\documentclass[sigconf,nonacm]{acmart}
\usepackage{enumitem}
\usepackage{booktabs}
\usepackage{amsmath}
\usepackage{amssymb}
\usepackage[normalem]{ulem}
\newcommand{\remove}[1]{\sout{#1}}
% ---- draft conveniences: suppress ACM reference block / conference banner ----
\settopmatter{printacmref=false, printccs=false}
\setcopyright{none}

\AtBeginDocument{%
  \providecommand\BibTeX{{%
    \normalfont B\kern-0.5em{\scshape i\kern-0.25em b}\kern-0.8em\TeX}}}

\begin{document}

\title{Motion Encoding for Human Perceptual Similarity}

\author{Bennie Bendiksen Gutierrez}
\email{b.bendiksen001@umb.edu}
\affiliation{%
  \institution{University of Massachusetts Boston}
  \city{Boston}
  \state{Massachusetts}
  \country{USA}
}

\author{Funda Durupinar}
\affiliation{%
  \institution{University of Massachusetts Boston}
  \city{Boston}
  \state{Massachusetts}
  \country{USA}
}

\renewcommand{\shortauthors}{Bendiksen Gutierrez and Durupinar}

\begin{abstract}
We ask whether a learned representation of human motion can predict {perceptual}
similarity---how similar two movements appear to human observers---more faithfully than
geometric distance measures (Dynamic Time Warping and a quaternion-geodesic distance) that
dominate motion analysis. Using Laban-annotated motion-capture performances of three
action types (walking, pointing, picking) paired with a triplet-based human similarity-rating
dataset, we show that a self-supervised {masked motion predictor augmented with an
auxiliary pose-reconstruction objective (MAMP+pose)} produces motion embeddings whose
{raw, unsupervised} distances align with human perception more strongly than both
geometric baselines {and} a generic learned motion encoder
across all three actions, under a strict, repeated, leakage-clean nested cross-validation. To
our knowledge, this is the first single encoder to achieve this. A progressive ablation
isolates the contribution of each objective and motivates the final design: the reconstruction
and masked-prediction objectives prove complementary---each strong where the other is weak---and
composing them is the only design competitive on all three actions. We further show that the
human ratings, when used to fine-tune the encoder under a perceptual objective, further improve
alignment on all three actions. All claims are reported as fold-averaged correlations with
calibrated variability, against baselines re-evaluated on the identical folds.
\end{abstract}

\keywords{motion similarity, self-supervised learning, perceptual similarity, Laban Movement
Analysis, masked motion prediction}

\maketitle

% =====================================================================
\section{Introduction}

\subsection{Motion similarity and geometric distances as measures of similarity}
Quantifying how {similar} two human movements are is foundational to motion retrieval,
synthesis, style transfer, and the perceptual study of movement. The dominant approaches are
geometric distances computed on raw kinematic features:

\begin{itemize}
\item {Rotation geodesic distance}---the distance on the rotation group $SO(3)$ between
corresponding joint orientations, aggregated over joints and time. It respects the manifold
structure of rotations but assumes (after temporal alignment) frame-wise correspondence.
\item {Dynamic Time Warping (DTW)}---a non-linear temporal alignment that factors out
speed variation and emphasizes the {shape} and {progression} of a gesture. DTW is the
most widely adopted geometric framework for time-series motion comparison precisely because it
models the temporal structure to which human observers are sensitive.
\end{itemize}

Though strong, well-motivated baselines, the scientific question is whether a {learned}
motion representation can predict human perceptual judgments better than they do.

\subsection{Contributions}
\begin{enumerate}
\item A {single self-supervised motion encoder (MAMP+pose)} whose {raw,
unsupervised} embedding distances beat both geometric baselines (DTW, rotation geodesic)
{and a generic learned (reconstruction-only SSL) motion encoder} on human-perceptual
similarity across all three action types---including the expressively sparse pointing
action---when evaluated on motion clips held out of pretraining, under a repeated, leakage-clean
nested cross-validation with all baselines re-evaluated on the identical folds. The learned
baseline shows that {merely} learning a motion representation is insufficient as it does not
uniformly beat geometry; the composition is what wins all three. To our knowledge this is the
first single encoder to do so on this heterogeneous action set, and it requires {no
perceptual supervision and metric-learning stage}.
\item A {mechanistic, progressive-ablation account} of {why} the
masked-motion-plus-pose composition is the right encoder: reconstruction objectives capture
static configuration but not dynamics; masked motion prediction is the complement (a picking
specialist that collapses on walking and pointing); and only composing the two yields an encoder
competitive across walking, pointing, and picking. Each step's result motivates the next design
choice, evidence-first.
\item A demonstration that the human ratings are productively used by {fine-tuning the
encoder} under a perceptual objective derived from the comparison study---an {anchored
ranking loss whose margin is the human dissimilarity gap itself} (hyperparameter-free and
metric-aligned), which further improves held-out alignment on all three actions.
\item A {methodological contribution}: a repeated, leakage-clean nested cross-validation
protocol---with baselines re-evaluated on the identical folds---needed to make reliable,
reproducible perceptual-similarity claims on a small, high-variance human-rating set, and a
characterization of the single-split pitfalls it corrects.
\end{enumerate}

% =====================================================================
\section{Data corpus}

\subsection{Stimulus space: Laban effort tuples}
Stimuli are motion-capture performances parameterized by Laban Movement Analysis (LMA)
{effort tuples} $(e_1, e_2, e_3, e_4) \in \{-1, 0, +1\}^4$, encoding the four Laban
effort dimensions (Space, Time, Weight, Flow). Following established LMA semantics, the number
of non-zero efforts has perceptual meaning: 0 non-zero $=$ neutral (1 class); 2 non-zero $=$
{state} (24 classes); 3 non-zero $=$ {drive} (32 classes); single-effort and
full-effort tuples are rare in natural motion and are excluded.

The perceptual study uses the {neutral + state + drive} subset: {57 effort classes
per action} ($1$ neutral $+ 24$ states $+ 32$ drives), performed for three action
types---{walking, pointing, picking}---for a total of 171 effort-class stimuli. Each
performance additionally has a left/right-mirrored twin; the LMA states/drives clips excluded
from encoder pretraining (Section 2.3) comprise $57 \times 3 \times 2 = 342$ clips.

\subsection{Motion representation}
Each clip is converted to a per-frame feature of {34 joints $\times$ 6-D rotation
representation}: joint index 0 carries the root position/orientation, indices 1--33 carry 6-D
continuous joint rotations. The 6-D representation (Gram--Schmidt-recoverable to $SO(3)$) is
preferred over quaternions for its continuity and gradient stability.

Each of the 204 feature channels ($34 \times 6$) is {z-score standardized} using a
per-channel mean and standard deviation computed once over the entire training corpus; values
are then clipped to $\pm 10$. The same fixed statistics are applied at training and evaluation,
so no evaluation-clip information enters the normalization.

{Fixed-length windowing.} The encoder operates on a fixed {120-frame} window (four
seconds at 30\,fps). Each clip is mapped to this length deterministically: clips longer than 120
frames are truncated to their first 120 frames; shorter clips are {right-padded by
repeating their final frame}. This identical rule is applied in pretraining, in embedding
extraction, and in held-out encoding---there is no length-dependent branching between train and
test. (The geometric baselines of Section 5.2, by contrast, operate on the {raw, unpadded}
sequences.)

{Clip-length heterogeneity.} Walking and picking clips are long (median 73 and 71 frames;
up to 137), whereas pointing clips are markedly shorter (median 53 frames; max 100). Because
short clips are tail-padded, pointing clips contribute the largest fraction of repeated,
frozen-pose frames under the 120-frame window---a consequence of their brevity that connects to
their limited kinematic expressiveness (Section 3).

\subsection{Pretraining corpus (self-supervised encoder)}
The encoder is pretrained self-supervised on a broad, heterogeneous motion corpus of
{5{,}795 BVH clips} drawn from three sources: the LMA effort performances (433 clips,
including the eval actions); the Bandai-Namco Research Motion Dataset [Kobayashi et al., 2023]
(3{,}077 clips of professionally captured everyday actions in multiple expressive styles); and
the CMU Graphics Lab Motion Capture Database (2{,}285 general-purpose performances).

{A single common skeleton.} The three sources were originally captured with different
marker sets, joint hierarchies, frame rates, and file conventions. A substantial preprocessing
effort retargeted {all} sources onto one canonical skeleton---the {CMU joint
hierarchy}---and re-expressed every clip as a uniform BVH file at a common {30\,fps}. This
retargeting is what makes the corpus trainable by a single encoder: every clip yields the
identical {34-joint $\times$ 6-D rotation} per-frame feature (Section 2.2), regardless of
source.

{Held-out construction.} To guarantee a {strict held-out evaluation}, all 342 LMA
neutral/state/drive clips (mirror twins included)---exactly the clips on which perceptual
similarity is evaluated---are excluded from pretraining (and from the encoder's own validation),
leaving $\sim$5{,}453 training clips. The encoder therefore never observes any evaluation
stimulus during representation learning. For the corpus-fair comparison against TMR (Section
6.3), we additionally train a {corpus-matched} MAMP+pose on a $14{,}143$-clip subset of AMASS
(spanning KIT, BioMotionLab, CMU, ACCAD, MPI and other AMASS constituents, retargeted to our
common skeleton). This matches the budget of the $\sim$14.6k-motion HumanML3D corpus on which TMR
is trained. HumanML3D is itself drawn predominantly from AMASS ($\approx$92\%, the remainder from
HumanAct12), so the two encoders are trained at a matched budget from a largely shared motion
source; our subset is AMASS-only and does not include the small HumanAct12 portion. This corpus
is likewise leakage-clean: it contains none of the LMA evaluation clips.

\subsection{Human perceptual ratings (ground truth)}
Similarity ratings were collected via {triplet (2-of-3) comparison trials}. In each trial
a participant viewed three motions---{Left (0)}, {Neutral (1)}, {Right
(2)}---and selected the two judged most similar. The dataset comprises {1{,}540 triplet
trials per action} (one per distinct non-neutral effort pair, with the neutral as the third
item), each rated by a {median of 11 participants} (range 10--18); presentation order was
randomized. The three pairings $\{(0,1),(0,2),(1,2)\}$ of a given triplet received normalized
choice frequencies summing to 1, interpretable as probabilities of selecting each pairing as
most-similar.

\subsubsection{Notation for the three choice frequencies}
Fix a triplet with {Left (0)} and {Right (2)} the two effort exemplars and
{Neutral (1)} the neutral exemplar. Write the three aggregated choice frequencies (each in
$[0,1]$, summing to 1) as $c(0,2)$ (Left--Right chosen), $c(0,1)$ (Left--Neutral), and $c(1,2)$
(Neutral--Right). $c(0,2)$ is the quantity of primary interest: how often observers judged the
two effort exemplars {mutually most similar}, an absolute measure of the perceived
similarity of that effort pair on a common $[0,1]$ scale.

\subsubsection{The evaluation target: a perceptual dissimilarity in $[0,1]$}
For each non-neutral effort pair we define the perceptual {dissimilarity} as
\[
    d_{\text{perc}} = 1 - c(0,2) \in [0,1].
\]
A pair almost always grouped together ($c(0,2)\to 1$) has $d_{\text{perc}}\to 0$; a pair almost
never grouped together has $d_{\text{perc}}\to 1$. {$d_{\text{perc}}$ is the single
ground-truth quantity against which every distance measure---the learned embedding distances and
both geometric baselines---is correlated} (Section 5.2). It uses only the absolute Left--Right
frequency and no triplet-internal contrasts, so it is method-agnostic.

\subsubsection{Perceptual training signals}
The perceptual fine-tuning of Section 6.4 requires a per-pair supervisory signal derived from the
triplets. The {reported method trains on the ordinal $d_{\text{perc}}$ relation} via a
ranking loss; a finer, neutral-referenced construction---the {dynamic alpha}---is one of
the tested alternatives, defined here. Per non-neutral effort pair we generate two directed
Left--Right alphas contrasting the direct Left--Right frequency against each neutral-mediated
alternative:
\[
    \alpha(0\!\to\!2) = c(0,2) - c(0,1), \qquad
    \alpha(2\!\to\!0) = c(0,2) - c(1,2).
\]
Each alpha lies in $[-1,+1]$; a positive alpha is evidence the pair should sit {close} in
embedding space, a negative alpha that it should sit {far}. These alphas are used only by
the neutral-anchored {alternative} arms, which Section 6.4 finds do not generalize. A note
on hygiene: the reported ranking method supervises on the {ordinal relation} induced by
$d_{\text{perc}}$, while evaluation scores the Spearman correlation between embedding distances
and $d_{\text{perc}}$ on the held-out test fold; because train, selection, and test folds are
disjoint (Section 5.2), no class's rating is ever both a training target and its own test.

\subsection{Metric of merit}
The ground truth is {ordinal} by construction (a ranking induced by choice frequencies)
and the learned metric optimizes ordinal triplet inequalities; {Spearman rank correlation
is the primary metric.} Pearson is reported as a magnitude-alignment supplement.

% =====================================================================
\section{Per-action motion structure}
The three actions differ markedly in {kinematic expressiveness}---how much of the body
moves, over how large a range, and over how long. We quantify this from the motion data,
averaging per-effort-class statistics over the 57 classes of each action ($28$ articulated
joints; a joint counts as \emph{active} in a clip when its across-frame variability exceeds a
small fixed threshold).

\begin{table}[h]
\centering
\caption{Per-action kinematic descriptors (mean over the 57 effort classes). Active-joint
fraction and temporal variability are reported as approximate, threshold-robust ranges; median
clip length is exact.}
\begin{tabular}{lccc}
\toprule
Descriptor & Walking & Pointing & Picking \\
\midrule
Active-joint fraction (of 28) & $\sim$0.9 & $\mathbf{\sim}$\textbf{0.55} & $\sim$0.7 \\
Relative temporal variability & high & \textbf{$\approx$ half} & high \\
Median clip length (frames)   & 73 & \textbf{53} & 71 \\
\bottomrule
\end{tabular}
\end{table}

{Pointing is, by every measure, the least kinematically expressive action.} On average only
\emph{about half} of the $28$ joints move meaningfully in a pointing clip ($\sim$$55\%$), versus
the large majority in walking ($\sim$$90\%$) and picking ($\sim$$70\%$); pointing's per-joint
temporal variability is roughly {half} that of the other two actions, and its clips are the
shortest (median $53$ vs $71$--$73$ frames). These gaps are large and hold across reasonable
choices of the activity threshold.
This expressive sparsity leads us to two hypotheses for pointing's distinctive difficulty:
{(H1, signal density)} fewer varying joints give the encoder less structure from which to
separate pointing effort classes, so raw pointing geometry should be less discriminative; and
{(H2, temporal brevity)} shorter clips afford the masked-motion pretext task fewer dynamic
regions to model. Under both, the prediction is the same: pointing's {raw} encoder distances
will be the least discriminative of the three actions. We see this borne out (Section 6): pointing
is the lowest of MAMP+pose's three raw correlations and the highest-variance action. What is not
foreordained is whether that residual signal is nevertheless enough to beat the geometric
baselines; we find that with the right pretraining objective it is.

% =====================================================================
\section{Encoder design---a progressive ablation toward MAMP+pose}
The encoder design was not chosen a priori; it is the endpoint of a sequence of controlled
training runs, each motivated by the failure of the previous. We present that sequence because it
establishes the central {two-capability composition} finding---that perceptual similarity
across heterogeneous actions requires an encoder modeling {both} motion dynamics and static
pose configuration---and rules out the obvious single-objective alternatives.

We evaluate every encoder identically: pool its output to one vector per clip, take L2 distances
between clips, and correlate against $d_{\text{perc}}$ under the repeated, leakage-clean nested
cross-validation of Section 5.2 (all Section 4 numbers are fold-averaged Spearman, $n=25$).
{No encoder is evaluated at a hand-chosen training epoch.} Within each fold the training
epoch is an inner-loop hyperparameter: a candidate checkpoint is taken on a {uniform
100-epoch grid across the model's full trajectory}, the checkpoint maximizing raw perceptual
Spearman on the disjoint inner validation fold is chosen, and only then is it scored on the
untouched TEST fold. The grid is identical for every model, removing checkpoint-selection as a
researcher degree of freedom. Encoders differ along two axes we deliberately separate---the
{autoencoder backbone} (Section 4.1) and the {training objective} (Sections
4.2--4.3).

\subsection{Step 1---backbone: the MLD SkipTransformer over a plain transformer}
We consider two transformer-autoencoder backbones for the reconstruction encoders: a
{plain transformer autoencoder} (a standard encoder/decoder with a compressed latent
bottleneck) and the {MLD-style transformer autoencoder} of Motion Latent Diffusion [Chen
et al., CVPR 2023]---a transformer encoder/decoder with {U-Net-like long skip connections}
(``SkipTransformer''), symmetric long-range links from each early encoder layer to the matching
late decoder layer. The skips carry high-frequency per-frame detail forward that the plain
bottleneck discards. Holding the {objective} fixed (reconstruction-only) and varying
{only} the architecture:

\begin{table}[h]
\centering
\caption{Reconstruction-only encoders, raw Spearman ($n=25$, mean $\pm$ SE).}
\setlength{\tabcolsep}{1pt}
\small
\begin{tabular}{lccc}
\toprule
Encoder & Walking & Pointing & Picking \\
\midrule
Plain transformer AE          & $+0.476 \pm 0.031$ & $+0.213 \pm 0.034$ & $+0.491 \pm 0.029$ \\
{MLD SkipTransformer AE} & $\mathbf{+0.532 \pm 0.023}$ & $\mathbf{+0.331 \pm 0.039}$ & $\mathbf{+0.498 \pm 0.029}$ \\
\bottomrule
\end{tabular}
\end{table}

The MLD SkipTransformer wins on {all three actions at the same objective}, decisively on
pointing ($+0.118$), consistent with the skips preserving high-frequency pose detail that the
plain bottleneck discards.

{Deterministic vs.\ variational.} We use the {deterministic} MLD autoencoder rather
than its variational form (MLD-VAE) on principled grounds, not performance: a VAE adds a KL term
toward a prior (a generative design choice) that confounds the reconstruction probe, and our
headline encoder is itself deterministic, so a deterministic comparator is apples-to-apples.
Empirically the two are comparable on the perceptually decisive pointing axis ($+0.331$ AE vs.\
$+0.324$ VAE); MLD-VAE achieves the strongest walking score in the sweep ($+0.597$) but remains a
walking/picking specialist that stalls on pointing.

\subsection{Step 2---objective: a controlled velocity-vs-reconstruction factorial}
Having fixed the backbone, we ask whether an explicit {velocity} term endows the encoder
with dynamics. To avoid arbitrary loss weighting, we run a {controlled single-variable
factorial}: each loss term is on (unit weight) or off, on {both} backbones.

\begin{table}[h]
\centering
\setlength{\tabcolsep}{1pt}
\small
\caption{Velocity factorial, raw Spearman ($n=25$, mean $\pm$ SE).}
\begin{tabular}{lccc}
\toprule
Encoder (objective) & Walking & Pointing & Picking \\
\midrule
\multicolumn{4}{l}{{Plain transformer}} \\
\quad reconstruction-only & $+0.476 \pm 0.031$ & $+0.213 \pm 0.034$ & $+0.491 \pm 0.029$ \\
\quad reconstruction $+$ vel & $+0.472 \pm 0.026$ & $+0.205 \pm 0.039$ & $+0.399 \pm 0.037$ \\
\quad velocity-only & $+0.321 \pm 0.033$ & $+0.362 \pm 0.036$ & $+0.471 \pm 0.033$ \\
\multicolumn{4}{l}{{MLD SkipTransformer}} \\
\quad reconstruction-only & $+0.532 \pm 0.023$ & $+0.331 \pm 0.039$ & $+0.498 \pm 0.029$ \\
\quad reconstruction $+$ vel & $+0.537 \pm 0.022$ & $+0.339 \pm 0.035$ & $+0.493 \pm 0.030$ \\
\quad velocity-only & $+0.451 \pm 0.031$ & $+0.278 \pm 0.031$ & $+0.254 \pm 0.039$ \\
\bottomrule
\end{tabular}
\end{table}

Three results. {(1) Adding a velocity term to reconstruction changes almost nothing}---flat
within noise on MLD, flat-to-worse on the plain transformer; on neither backbone does it move
pointing across the $+0.370$ bar. {(2) Velocity-only is a weaker objective overall and
reshapes {which} action is served}---with reconstruction removed, global quality drops (MLD
picking collapses $+0.498 \to +0.254$). {(3) The telling signal:} velocity-only is the
only reconstruction-family cell that lifts plain-transformer pointing, and it lifts it markedly
($+0.213 \to +0.362$), while degrading walking. We are careful about what this licenses---it does
not tell us where pointing's perceptual signal physically resides---but it does establish that a
purely temporal target carries perceptual signal a pose target does not surface, a first
indication that {dynamics information is perceptually relevant} and under-exploited by
reconstruction. Yet velocity as an added {loss} cannot capitalize on it (result 1). The
lesson is therefore not ``add velocity'' but ``capture dynamics properly,'' motivating the move
to a task that {forces} dynamics into the representation.

\subsection{Step 3---why a velocity loss cannot supply dynamics}
A velocity loss on a visible reconstruction is {partly redundant}. A reconstruction
autoencoder sees the entire clip and is rewarded for reproducing every frame; if it reconstructs
the poses $x_t$ accurately, the frame-to-frame differences $x_{t+1}-x_t$---the velocities---are
reproduced automatically. A velocity term is therefore largely already satisfied and supplies
little additional gradient. This predicts precisely what the factorial shows. The lesson that
carries forward: {dynamics must be forced into the representation by the learning task
itself}---by making the model {predict} motion it cannot see---not appended as a loss term
on a fully-visible reconstruction.

\subsection{Step 4---masked motion {prediction}, not reconstruction}
The resolution comes from the Masked Motion Prediction (MAMP) framework [Mao et al., ICCV 2023].
Rather than reconstructing visible poses, MAMP {masks $\sim$80\% of spatio-temporal
patches and predicts the {motion} (temporal deltas) of the masked regions} from sparse
visible context, with a motion-magnitude prior biasing masking toward dynamically rich regions.
This imposes genuine dynamics-modeling pressure: the encoder cannot copy what it cannot see.

{MAMP alone is a dynamics specialist.} The motion-only MAMP encoder's single strong action
is picking ($+0.490$ raw, the one bar it clears), while it is weak on walking ($+0.351$) and
pointing ($+0.299$). What is solid---and load-bearing---is the {empirical complementarity}
between the two objectives: masked-motion prediction succeeds on exactly the action (picking)
where reconstruction was weakest, and fails on exactly the actions (walking, pointing) where
reconstruction was strongest. Neither objective alone serves all three, which motivates
{composing} them.

\subsubsection{Pretraining budget and the pretext-vs-downstream divergence}
We adopt the original MAMP pretraining hyperparameters and pretrain to convergence of the
self-supervised loss, saving the full checkpoint trajectory; the reported epoch is selected per
fold on the inner validation fold over the uniform 100-epoch grid. For the MAMP family this
consistently selects an {early} checkpoint---the held-out perceptual signal rises quickly
and plateaus by a few hundred epochs while the self-supervised loss keeps descending---a concrete
instance of the pretext-vs-downstream divergence. Selection is leakage-clean (the validation fold picks;
TEST fold reports).

\subsection{Step 5---the pose augmentation (MAMP+pose)}
Masked motion prediction supplies dynamics but discards the static-configuration fidelity that
reconstruction recovered---and with it walking and pointing. We augment it with an
{auxiliary pose-reconstruction objective}. In standard MAMP a single linear head maps each
masked patch's decoder feature to its {motion} target (the temporal delta $x_{t+s}-x_t$);
we add a {second, parallel linear head} off the {same} decoder features mapping each
masked patch to its {pose} target (the raw per-frame value $x_t$). The training loss is
$L = L_{\text{motion}} + \lambda_{\text{pose}} \cdot L_{\text{pose}}$ with
$\lambda_{\text{pose}}=1.0$. Because both heads read the same decoder features and backpropagate
into the same encoder, the encoder is pressured to produce one latent that simultaneously encodes
how the body is {moving} (the dynamics that win picking) and how it is {configured} (the
static fidelity walking and pointing depend on). Two parallel heads (rather than summing targets)
let each target be predicted in its natural form---a temporal delta and an absolute pose live on
different scales---while sharing all representation-bearing layers. Section 6.5 shows this
composition rescues walking and pointing while retaining picking.

% =====================================================================
\section{The perceptual inference pipeline}

\subsection{Raw-embedding regime (unsupervised)}
A clip is mapped to the fixed 120-frame window, standardized, and patchified into spatio-temporal
patches---one patch per (4-frame block $\times$ joint), giving $30 \times 34 = 1020$ patches. The
pretrained encoder (run with mask ratio 0, all patches visible) produces a 256-D feature per
patch; we mean-pool to a single {256-D clip embedding} (the encoder-native linear-probe
readout). Pairwise {Euclidean (L2) distances} between clip embeddings are correlated
(Spearman) against $d_{\text{perc}}$ over the held-out test subset. This regime uses {no
human ratings} during representation learning and therefore measures the encoder's intrinsic
perceptual alignment.

\subsection{Evaluation protocol and baselines}
{Repeated, leakage-clean nested cross-validation.} All reported correlations are the mean
over {5 repeats $\times$ 5 outer folds ($n=25$)}. Effort classes are partitioned
(magnitude-stratified) into 5 folds with a repeat-specific seed. For each outer fold: the fold is
the {test} set (reported, never seen during training or selection); one disjoint fold is
the {validation} fold (the model-selection set), used {only} to pick each encoder's checkpoint
by raw held-out Spearman; the remaining folds are the {training} set. The validation fold plays
the standard role of a held-out validation set within each outer fold: all model-selection
decisions read it, and it is disjoint from the reported test fold. The encoder additionally never saw
these clips during pretraining, so the test fold is held out of both stages. We report mean $\pm$
standard error (SE $=$ SD$/\sqrt{25}$).

{Why not a single fixed split.} The per-action correlation is strongly split-dependent,
most severely for pointing, whose value ranges roughly $[0.23, 0.55]$ depending only on which
classes land in the held-out subset. Conclusions drawn from one split do not survive
re-partitioning; the repeated nested-CV averages over this variance and is the basis for every
claim.

{Baselines on identical folds.} Both geometric baselines are computed on the {raw,
unpadded} motion features over the {same 25 test folds} and correlated against the same
$d_{\text{perc}}$. A representational note: the geometric baselines operate on the native
motion-capture rotation representation---per joint a unit \emph{quaternion} ($28$ joints
$\times\,4 = 112$ channels per frame)---whereas the learned encoder consumes the 6-D continuous
rotation features of Section~2.2; each representation is the natural one for its method, and both
describe the same underlying joint rotations.

Both baselines use Dynamic Time Warping for temporal alignment but differ in the per-frame local
cost. The \textbf{DTW} baseline uses a \textbf{Euclidean} per-frame cost between the (flattened)
quaternion feature vectors, and reports the length-normalized accumulated cost of the optimal
monotone alignment---so it compares the two motions' shape and progression while absorbing
differences in speed and duration. The \textbf{geodesic} baseline replaces the Euclidean local
cost with a \textbf{quaternion geodesic}: for two frames with per-joint unit quaternions $q_1, q_2$,
the frame cost is the mean over joints of $2\arccos(|\langle q_1, q_2\rangle|)$ (the geodesic angle
on the rotation manifold, invariant to the quaternion double cover); DTW aligns the sequences under
this cost and we report the length-normalized aligned mean. ({Correction to a legacy
implementation}: an earlier version padded clips to a fixed length by repeating the final frame and
compared frame-wise without alignment; for short pointing clips this injected long runs of
frozen-pose frames, spuriously deflating distances. Computing on raw, unpadded sequences with
proper temporal alignment removes this bias and \emph{strengthens} the baseline.)

% =====================================================================
\section{Results}

\subsection{Headline---held-out, repeated nested cross-validation}
We report Spearman correlation between each method's pair distances and human dissimilarity
$d_{\text{perc}}$. All headline numbers are the mean over the repeated, leakage-clean nested
cross-validation; the geometric baselines are computed on the identical 25 folds and fold-averaged
the same way.

\begin{table}[h]
\centering
\setlength{\tabcolsep}{1pt}
\small
\caption{Raw MAMP+pose vs.\ geometric and learned baselines ($n=25$, mean $\pm$ SE).}
\begin{tabular}{lccc}
\toprule
Method & Walking & Pointing & Picking \\
\midrule
{MAMP+pose, raw} & $\mathbf{+0.552 \pm 0.026}$ & $\mathbf{+0.434 \pm 0.030}$ & $\mathbf{+0.543 \pm 0.033}$ \\
DTW (geometric)        & $+0.490 \pm 0.028$ & $+0.273 \pm 0.039$ & $+0.303 \pm 0.038$ \\
Geodesic (geometric)   & $+0.371 \pm 0.028$ & $+0.234 \pm 0.031$ & $+0.264 \pm 0.030$ \\
Vanilla SSL (recon-only) & $+0.476 \pm 0.031$ & $+0.213 \pm 0.034$ & $+0.491 \pm 0.029$ \\
\bottomrule
\end{tabular}
\end{table}

On Spearman, the {raw, unsupervised MAMP+pose encoder beats both geometric baselines and
the learned baseline on all three actions} under matched held-out evaluation---to our knowledge
the first single encoder to do so. Treating each fold-difference as paired across the 25 matched
folds, the pointing and picking advantages over the geometric baselines are large and
well-separated from zero ($\approx$3--5$\sigma$ on a conservative independent-SEM test); the one
near-marginal cell is walking vs.\ DTW ($+0.062$, $\approx$1.6$\sigma$), though walking clears
geodesic by a wide margin. The all-three result is a property of the raw encoder geometry and does
not require any learned metric refinement.

\subsection{The raw-encoder result---a stronger, unsupervised claim}
The all-three win is achieved by the {raw} encoder embeddings---no perceptual supervision
whatsoever. The perceptual-similarity structure is therefore {intrinsic to the
self-supervised representation}, not imposed by a downstream metric. Beating a heavily-optimized
temporal baseline like DTW with unsupervised embeddings---including on the expressively sparse
pointing action, where DTW's time-warping was expected to dominate---indicates the masked-motion
pretext task, augmented with the pose head, genuinely captures the kinematic structure underlying
human similarity judgments.

\subsection{Comparison to a supervised learned representation (TMR)}
A sterner test is a {purpose-built, supervised} motion representation. We compare against
{TMR} [Petrovich et al., ICCV 2023], a text$\leftrightarrow$motion retrieval model trained
with contrastive language supervision on HumanML3D---a representation shaped by a {large}
corpus of semantic annotations, the opposite supervision regime to ours.

{Protocol and a conservative handicap.} TMR consumes HumanML3D 263-D motion features, so we
convert each rated clip from BVH to SMPL parameters (joints2smpl SMPLify fit, mean reconstruction
error $\approx$4\,cm $\approx$2\% of body height) and then to the 263-D features TMR expects,
encode through the released HumanML3D TMR motion encoder, and evaluate on the {identical 25
nested-CV folds} and the same $d_{\text{perc}}$. The $\approx$4\,cm conversion perturbs TMR's input
slightly, so any MAMP+pose advantage is a {conservative lower bound}. For a corpus-fair
comparison we use the corpus-matched MAMP+pose (Section 2.3) as the comparator.

\begin{table}[h]
\centering
\caption{Corpus-matched MAMP+pose vs.\ TMR ($n=25$, mean $\pm$ SE).}
\setlength{\tabcolsep}{1pt}
\small
\begin{tabular}{lccc}
\toprule
Method & Walking & Pointing & Picking \\
\midrule
TMR (supervised, HumanML3D) & $+0.452 \pm 0.025$ & $\mathbf{+0.425 \pm 0.036}$ & $+0.408 \pm 0.029$ \\
MAMP+pose (matched, raw)    & $+0.462 \pm 0.027$ & $+0.302 \pm 0.031$ & $+0.481 \pm 0.031$ \\
{\ + ranking fine-tune} & $\mathbf{+0.699 \pm 0.022}$ & $\mathbf{+0.463 \pm 0.035}$ & $\mathbf{+0.605 \pm 0.031}$ \\
\bottomrule
\end{tabular}
\end{table}

Two honest reads. {Raw}, the corpus-matched MAMP+pose beats TMR on walking and picking but
{TMR is stronger on pointing} ($+0.425$ vs.\ $+0.302$)---the one action where abundant
semantic supervision helps most, and where our raw encoder is weakest. TMR is the strongest
pointing baseline in the entire study (far above DTW's $+0.273$) and, unlike the geometric
baselines, does not collapse on pointing---it is roughly flat across actions, consistent with a
semantically supervised representation capturing broad motion structure uniformly but not tuned to
perceptual effort-similarity. {With the light ranking fine-tune (Section 6.4), MAMP+pose
beats TMR on all three actions}, including pointing ($+0.463$ vs.\ $+0.425$), achieved with a far
smaller supervision set ($\approx$1{,}540 within-task perceptual triplets per action) than TMR's
corpus of semantic annotations. The finding is one of {alignment over abundance}: sparse
supervision aligned to the target task matches or exceeds abundant supervision aligned to a
different task. We state the pointing result precisely---TMR wins pointing on the raw encoder and
is edged only after the perceptual fine-tune---rather than overclaiming a raw all-three win against
this baseline.

\subsection{Perceptual fine-tuning: using the human ratings to improve the encoder}
The raw encoder already wins. We now ask whether the small budget of human ratings can push it
further. The productive way, we find, is to back-propagate a perceptual objective {into the
encoder itself}, reshaping the representation so that embedding distances track human
dissimilarity---rather than fitting a separate metric on a frozen embedding.

{An anchored ranking loss with a preference-scaled margin.} The ground truth
$d_{\text{perc}}$ is {ordinal}, so the natural objective is a {ranking} loss. For an
anchor class $i$ and two others $j,k$, let $p_j = d_{\text{perc}}(i,j)$ and $p_k =
d_{\text{perc}}(i,k)$; the human ordering names a {near} and a {far} class. With pairwise
distances $\hat d$ min--max normalized to $[0,1]$ within the batch, we impose
\[
    L = \mathrm{relu}\big(\hat d(i,\text{near}) - \hat d(i,\text{far}) + m\big),
    \qquad m = |p_k - p_j|,
\]
i.e.\ the anchor--near distance must fall below the anchor--far distance by a margin equal to the
{human dissimilarity gap} itself. Two properties make this the principled choice: the
margin is {hyperparameter-free and adaptive} (confidently-ordered pairs are pushed apart
harder than near-ties, with no margin constant to tune), and the loss is {metric-aligned by
construction}---it optimizes the same ordinal quantity, over the same non-neutral pairs, that the
Spearman evaluation scores.

{Leakage-clean protocol and a single held-out selection.} Fine-tuning runs under the same
repeated nested-CV as the raw results; the only change is that within each fold's training portion
the MAMP+pose encoder is fine-tuned rather than frozen. Two design points make the reported numbers
conservative. First, the base checkpoint is chosen {per fold on the validation fold only}.
Second, we use a {single} held-out selection: rather than also early-stopping on the
validation fold, we train a {fixed budget and report the final epoch}---removing any
``selection used twice'' concern by construction.

\begin{table}[h]
\centering
\caption{Ranking fine-tune vs.\ raw MAMP+pose ($n=25$, mean $\pm$ SE).}
\setlength{\tabcolsep}{1pt}
\small
\begin{tabular}{lccc}
\toprule
Method & Walking & Pointing & Picking \\
\midrule
raw MAMP+pose (no ratings) & $+0.462 \pm 0.027$ & $+0.302 \pm 0.031$ & $+0.481 \pm 0.031$ \\
{\ + ranking fine-tune} & $\mathbf{+0.699 \pm 0.022}$ & $\mathbf{+0.463 \pm 0.035}$ & $\mathbf{+0.605 \pm 0.031}$ \\
\bottomrule
\end{tabular}
\end{table}

The ranking fine-tune {lifts all three actions substantially} from the corpus-matched raw
base. It is robust to the stopping rule: reporting the fixed-budget final epoch matches---in fact
slightly exceeds---a variant that additionally early-stops on the validation fold, so the
single-selection design loses nothing. This follows directly from the fine-tuning dynamics,
which we profile across all $25$ folds: the mean validation-fold Spearman rises steeply from
$0.44$ at initialization and \textbf{plateaus by roughly epoch $100$--$120$} at $\approx 0.60$,
then holds---drifting slightly downward toward the $200$-epoch budget---while the \emph{training}
loss continues to descend monotonically. The fixed $200$-epoch budget therefore lands the encoder
squarely on this converged plateau: past the point of diminishing validation return, but before
meaningful overfitting. Consequently the exact stopping epoch is immaterial and a second,
early-stopping use of the validation fold is unnecessary. The pattern mirrors the
pretext-versus-downstream divergence of self-supervised pretraining (Section~4.4.1)---the
optimized loss keeps improving after the perceptually-relevant signal has saturated. (A
direct-regression objective and a neutral-anchored,
alpha-margin objective were also tried; both under-perform the ranking loss, the latter failing to
generalize on pointing---consistent with the perceptual signal residing in the direct inter-state
ordering the ranking loss targets rather than in neutral-referenced context shifts.) In short, the
human ratings are best used not as a metric fit on a fixed representation but as a {light,
ranking-based fine-tuning of the encoder}.

\subsection{Ablation---the pose head is necessary}
\begin{table}[h]
\centering
\setlength{\tabcolsep}{1pt}
\small
\caption{Motion-only MAMP vs.\ MAMP+pose ($n=25$, mean $\pm$ SE).}
\begin{tabular}{lccc}
\toprule
Objective & Walking & Pointing & Picking \\
\midrule
motion-only (MAMP) & $+0.351 \pm 0.032$ & $+0.299 \pm 0.032$ & $+0.490 \pm 0.039$ \\
{\ + pose head (MAMP+pose)} & $\mathbf{+0.552 \pm 0.026}$ & $\mathbf{+0.434 \pm 0.030}$ & $\mathbf{+0.543 \pm 0.033}$ \\
\bottomrule
\end{tabular}
\end{table}

Motion-only masked prediction clears the bar only on picking, its best action, and is weak on
walking and pointing. Adding the pose head {lifts every action}, most sharply walking
($+0.351 \to +0.552$) and pointing ($+0.299 \to +0.434$, the gain that carries it clear of the
baselines), while retaining picking. This confirms the two-capability composition: the motion
target supplies dynamics, the pose target supplies configuration fidelity, and the combination is
the only encoder clearing both baselines on all three actions.

\subsection{Summary of the per-action story}
A single self-supervised encoder, with no perceptual supervision, beats both geometric baselines
on all three actions under repeated nested cross-validation---including the expressively sparse
pointing action. Where the human ratings help further, it is by fine-tuning the encoder (Section
6.4), which lifts all three actions above the already-baseline-beating raw encoder.

\subsection{Full model sweep (all encoders, raw)}
All values are from the same repeated nested cross-validation ($n=25$), evaluated on the canonical
57 effort classes per action; each encoder's checkpoint is selected per fold on the disjoint
validation fold from an 11-checkpoint uniform 100-epoch grid. The geometric bars each encoder must
clear (the stronger baseline per action) are {walking $>0.478$, pointing $>0.370$, picking
$>0.431$}; we mark $\checkmark$ when the encoder's 95\% CI lower bound (mean $-1.96\cdot$SE) clears
that bar.

\begin{table}[h]
\centering
\caption{Raw embeddings, Spearman ($n=25$, mean $\pm$ SE).}
\setlength{\tabcolsep}{1pt}
\small
\small
\begin{tabular}{lcccc}
\toprule
Encoder (backbone, objective) & Walking & Pointing & Picking & all 3? \\
\midrule
Plain, recon-only    & $+0.476{\pm}0.031$ & $+0.213{\pm}0.034$ & $+0.491{\pm}0.029$ & --- \\
Plain, recon$+$vel   & $+0.472{\pm}0.026$ & $+0.205{\pm}0.039$ & $+0.399{\pm}0.037$ & --- \\
Plain, velocity-only & $+0.321{\pm}0.033$ & $+0.362{\pm}0.036$ & $+0.471{\pm}0.033$ & --- \\
MLD AE, recon-only   & $+0.532{\pm}0.023$ & $+0.331{\pm}0.039$ & $+0.498{\pm}0.029$ & --- \\
MLD AE, recon$+$vel  & $+0.537{\pm}0.022$ & $+0.339{\pm}0.035$ & $+0.493{\pm}0.030$ & --- \\
MLD AE, velocity-only& $+0.451{\pm}0.031$ & $+0.278{\pm}0.031$ & $+0.254{\pm}0.039$ & --- \\
MLD VAE, recon-only  & $+0.597{\pm}0.021$ & $+0.324{\pm}0.032$ & $+0.501{\pm}0.029$ & --- \\
MAMP (motion-only)   & $+0.351{\pm}0.032$ & $+0.299{\pm}0.032$ & $+0.490{\pm}0.039$ & --- \\
{MAMP+pose}   & $\mathbf{+0.552{\pm}0.026}$ & $\mathbf{+0.434{\pm}0.030}$ & $\mathbf{+0.543{\pm}0.033}$ & $\checkmark$ \\
\bottomrule
\end{tabular}
\end{table}

{MAMP+pose is the only encoder whose 95\% CI clears all three baseline bars.} Pointing is
the discriminator: the strongest competitors (the MLD AEs) reach walking/picking comfortably but
stall on pointing at $+0.32$--$0.34$ (CI not clearing $0.370$), exactly the expressively-sparse
action the composition is built to capture. No competing encoder is strong on all three actions;
each alternative is a specialist. The architecture and objective axes are separable and both
matter: at a fixed reconstruction objective the MLD backbone beats the plain transformer, while at
a fixed backbone it is the masked-motion-plus-pose objective---not reconstruction with or without a
velocity term---that produces the only all-three winner. Section 6.8 tightens this into a direct
de-confound.

\subsection{Disentangling architecture from objective---is it the pose head or the backbone?}
The step from the reconstruction encoders to MAMP changes {two} things at once---the
{objective} (reconstruction $\to$ masked prediction) and, implicitly, the {backbone},
since the MLD reconstruction encoders carry U-Net skip connections while the MAMP encoder is a
plain (skip-free) transformer. A skeptic could argue that MAMP+pose recovers configuration not
because of the pose {objective} but because skip connections---the mechanism that lifts
reconstruction pointing from $+0.213$ to $+0.331$---would do the same for masked prediction. We
test this directly.

{An MLD-style encoder inside MAMP.} We port the MLD SkipTransformer's {encoder-internal}
U-Net into the MAMP encoder, holding the masked-motion objective, corpus, normalization, and split
fixed. Because MAMP masks $\approx$80\% of tokens before the encoder, we run the U-Net encoder on
the {full, unmasked} sequence and route only the visible-position features to the decoder,
so the encoder is train/inference-consistent and its skips operate on complete motion while the
prediction task remains non-trivial (masked targets never leak to the decoder). We call this
{MAMP-uencfull}; crucially it has {no pose head}, isolating the effect of the
MLD-style {architecture} under the masked objective.

\begin{table}[h]
\centering
\caption{Architecture vs.\ objective on the masked-prediction backbone ($n=25$, mean $\pm$ SE).}
\small
\setlength{\tabcolsep}{1pt}

\begin{tabular}{lcccc}
\toprule
Encoder & Walking & Pointing & Picking & all 3? \\
\midrule
MAMP (plain enc, motion-only)     & $+0.351{\pm}0.032$ & $+0.299{\pm}0.032$ & $+0.490{\pm}0.039$ & --- \\
MAMP-uencfull (U-Net enc, motion-only) & $+0.470{\pm}0.029$ & $+0.275{\pm}0.028$ & $+0.566{\pm}0.030$ & --- \\
{MAMP+pose (plain enc, $+$pose)} & $\mathbf{+0.552{\pm}0.026}$ & $\mathbf{+0.434{\pm}0.030}$ & $\mathbf{+0.543{\pm}0.033}$ & $\checkmark$ \\
\bottomrule
\end{tabular}
\end{table}

The result is unambiguous on the discriminating axis. Giving MAMP the MLD-style U-Net encoder
{does not recover pointing}: it moves from $+0.299$ to $+0.275$---statistically flat, if
anything slightly lower---and remains below even the MLD reconstruction encoders' pointing and far
below the bar. The pose {objective}, by contrast, lifts pointing to $+0.434$ ($\approx$4\,SE
gain). {The configuration recovery that carries pointing is supplied by the
pose-reconstruction objective, not by the encoder architecture.} Two secondary observations sharpen
the mechanism: the U-Net encoder is not inert---it lifts the coarse-motion actions (walking $+0.351
\to +0.470$, and the best picking in the sweep, $+0.566$)---it simply does not manufacture the fine
configuration signal the sparse pointing action requires; and we run it on the full, unmasked
sequence deliberately, as the architecture's best case, since skips taken after masking (on the
visible $\approx$20\% only) are uniformly weaker.

We complete the $2\times 2$ of architecture $\times$ objective, since MLD's own recipe couples both
skip connections and a reconstruction objective. The remaining cell---an MLD-style U-Net encoder
trained {with} the pose objective---asks whether the two compound.

\begin{table}[h]
\centering
\setlength{\tabcolsep}{1pt}
\small
\caption{Architecture $\times$ objective (raw Spearman, $n=25$, mean $\pm$ SE).}
\begin{tabular}{lcc}
\toprule
 & plain encoder & MLD-style U-Net encoder \\
\midrule
motion-only & MAMP: $+0.299{\pm}0.032$ & uencfull: $+0.275{\pm}0.028$ \\
$+$ pose    & MAMP+pose: $\mathbf{+0.434{\pm}0.030}$ & {[pending]} \\
\bottomrule
\end{tabular}
\end{table}

{[Placeholder---the MAMP-uencfull+pose cell is pending a training run; to be filled with the
$n=25$ mean $\pm$ SE on completion. If its pointing exceeds $+0.434$, architecture and objective are
complementary---the U-Net encoder's benefit unlocks only with the pose objective, echoing MLD's own
coupling of skips with reconstruction; if $\approx +0.434$, the pose objective saturates
configuration recovery and the architecture adds nothing on top. Either outcome does not alter the
conclusion that the pose objective, not the backbone, recovers configuration.]}

% =====================================================================
\section{Why MAMP+pose succeeds: modeling vs.\ copying}
The central mechanistic claim is the distinction between {modeling} dynamics and
{copying} poses. Reconstruction objectives present the model with the entire visible motion
and reward reproducing it; under full visibility, the velocities are reproduced for free as a
byproduct of accurate pose copying, so a velocity term is partly redundant. Velocity-penalized
reconstruction thus remains a {pose-biased} representation. Masked motion prediction removes
the copying shortcut: with $\sim$80\% of patches hidden and the motion of the hidden, high-energy
regions as the target, the encoder must infer kinematic progression it cannot observe---it must
internalize how the body moves, not merely where it currently is. This is why motion-only masked
prediction is the empirical complement of reconstruction. The auxiliary pose head restores the
static-configuration fidelity reconstruction had. The {composition}---masked dynamics
modeling plus pose reconstruction---supplies both capabilities in one encoder and is the only
configuration competitive across all three heterogeneous actions; neither objective alone is.

% =====================================================================
\section{Reproducibility}

\subsection{Evaluation protocol (exact)}
{Repeated nested cross-validation.} 5 repeats $\times$ 5 outer folds ($n=25$).
Magnitude-stratified folds per action with seed $42 + 1000\cdot\text{repeat}$. Per outer fold $f$:
test $=$ fold $f$; selection $=$ fold $(f{+}1)\bmod 5$; training $=$ remaining three folds. The
neutral exemplar is added to every subset. Checkpoint selection uses raw held-out Spearman on the
validation fold (never the test fold). For the perceptual fine-tuning the validation fold is used
{once}, to pick the per-fold base checkpoint; the fine-tune itself runs a {fixed epoch
budget} (no second, early-stopping use of the validation fold).

{Evaluation universe.} The canonical 57 effort classes per action. The stored embedding
directories also contain non-canonical clips (single-effort and fully-polarized tuples) carrying
{no} human ratings; these cannot enter any correlation, but if left in the pool that is
shuffled to form the split they perturb which rated classes land in each fold and silently change
every number---most severely for pointing. The split must shuffle the canonical 57 only. (An
earlier inability to reproduce the pointing baseline traced entirely to this.)

{Per-pair target.} $d_{\text{perc}} = 1 - \text{count\_normalized}[(\text{selected}_0{=}0,
\text{selected}_1{=}2)]$. Correlation is Spearman of pooled-embedding L2 distance vs.\
$d_{\text{perc}}$ over exactly the test-fold pairs that carry a human rating, using plain pairwise
L2 with no neutral-centering.

{Strict hold-out and single corpus.} The 342 LMA evaluation clips are excluded from
pretraining and its validation; all reported encoders are trained on the same 3-dataset corpus, so
cross-encoder comparisons are corpus-matched.

\subsection{Pitfalls (for faithful reproduction)}
\begin{itemize}
\item {Single-split point estimates are unreliable.} Pointing-raw Spearman has fold-to-fold
SD $\approx 0.15$; earlier single-split numbers were partition artifacts. Use the repeated nested-CV
mean, reported as mean $\pm$ SE.
\item {Checkpoint identity.} The same encoder appears under two pretraining budgets in our
archive; always confirm the checkpoint epoch (and md5) before comparing embeddings.
\item {Code/version parity.} Reproduction requires the exact metric MLP, embedding loader,
and corrected geodesic function present on the execution host; stale copies change the output.
\end{itemize}

\subsection{Artifacts and policy}
Encoder checkpoints (with md5s), the LMA data feeder, the embedding extractor, the nested-CV and
baseline harnesses, the perceptual-fine-tune trainer, and the evaluation scripts are archived for
end-to-end reproduction, along with the exact per-fold random seeds. The corrected geodesic baseline
is computed on raw unpadded sequences (no frame-padding bias). Epoch counts are matched within a
backbone family but not forced across families, where an epoch is not a comparable unit.

% =====================================================================
\section{Limitations and scope}
\begin{itemize}
\item {Pointing's narrower margin and higher variance.} MAMP+pose's raw encoder does beat
both geometric baselines on pointing, but by the smallest margin and with the highest fold-to-fold
variance; against the supervised TMR baseline the raw encoder is edged on pointing and wins only
after the perceptual fine-tune.
\item {Three action types.} Generalization to a broader action taxonomy is untested; the
per-action heterogeneity we exploit may present differently elsewhere.
\item {Rating-set size.} The human rating data ($\sim$1{,}540 trials/action) is modest,
bounding how much the perceptual fine-tuning can reshape the encoder; a larger set might yield
larger or more stable gains.
\item {Polarized effort space.} We polarize each bipolar effort factor to $\{-1,0,+1\}$, so
$d_{\text{perc}}$ measures dissimilarity among polar-effort combinations, not along the continuous
intensity axis; graded-intensity perceptual distance is future work.
\item {Learned-baseline conversion.} The TMR comparison requires a BVH$\to$SMPL$\to$263-D
conversion ($\approx$4\,cm fit error) that slightly perturbs TMR's input; we frame our result as a
conservative lower bound accordingly.
\end{itemize}

% ---- plain-text citations retained for now; formalize into references.bib later ----
% [Chen et al., CVPR 2023] Motion Latent Diffusion (MLD SkipTransformer)
% [Mao et al., ICCV 2023] Masked Motion Prediction (MAMP)
% [Petrovich et al., ICCV 2023] TMR: text-to-motion retrieval
% [Kobayashi et al., 2023] Bandai-Namco Research Motion Dataset
% CMU Graphics Lab Motion Capture Database (http://mocap.cs.cmu.edu/)

\end{document}
